\section{Attraction of every orbit}\label{sec:orbits}

The registers in Section~\ref{sec:construction} have two tasks: they erase
all previous data before each preparation interval, and they supply additions
with the same phase as the main square at each subsequent readout.  We now
prove that these additions force every orbit whose lifted clock is unbounded
to escape to infinity.  A bounded clock instead forces the main coordinate
to zero.

Write $\mathbf v=(v_1,\ldots,v_q)$ and define the compact sets
\begin{equation}\label{eq:attracting-sets}
\begin{split}
 Z_0&=\{(50\bmod100,0,\mathbf v):\mathbf v\in\Sigma^q\},\\
 Z_\infty&=\{(t,\infty,0,\ldots,0):t\in\mathbb R/100\mathbb Z\}.
\end{split}
\end{equation}
Both are forward invariant under~\eqref{eq:map}.  On $Z_0$ one has
$h(50,0)=0$, all readout coefficients and inputs vanish, and the registers
update by $v_\ell\mapsto LS(v_\ell)$.  On $Z_\infty$ the main coordinate stays at
infinity, the clock advances by one, and the zero registers remain zero.

We keep the original time indices throughout the proof.  Thus, for an
arbitrary orbit, write
\[
 f^n(x)=(t_n\bmod100,z_n,v_{n,1},\ldots,v_{n,q}),\qquad
 t_{n+1}=t_n+h(t_n,z_n),
\]
where $t_0\in\mathbb R$ is any lift of the initial clock.  In particular,
\begin{equation}\label{eq:clock-increments}
 0\le t_{n+1}-t_n\le1.
\end{equation}

\begin{proposition}\label{prop:attraction}
For the map $f:M\to M$ constructed in Section~\ref{sec:construction}, every
$x\in M$ satisfies
\[
 \operatorname{dist}\bigl(f^n(x),Z_0\cup Z_\infty\bigr)\longrightarrow0
 \qquad(n\longrightarrow\infty)
\]
in every compatible metric on $M$.
More precisely, if the lifted clock of the orbit is bounded above, then
the clock converges to $50$ modulo $100$ and $z_n\to0$.
If that lift is unbounded above, then $z_n\to\infty$ in $\Sigma$ and all
registers are identically zero from some time onward.
\end{proposition}

\subsection{Bounded clocks and complete passages}

\begin{lemma}\label{lem:bounded-clock}
If $(t_n)$ is bounded above, then $t_n\to50+100k$ for some integer $k$ and
$z_n\to0$.  No condition on the initial register values is required.
\end{lemma}

\begin{proof}
Monotonicity gives a finite limit $t_*$, and hence
$h(t_n,z_n)=t_{n+1}-t_n\to0$.  Since $0\le b(t_n)\le h(t_n,z_n)$,
continuity and the zero set of $b$ imply $t_*=50+100k$.
For all sufficiently large $n$ the clock therefore avoids every readout
support, so $z_{n+1}=z_n^2$.  Moreover, $|z_n|\ge3/4$, including $z_n=\infty$,
would give $e(z_n)=1$ and $h(t_n,z_n)=1$.  Thus eventually $z_n$ is finite
with $|z_n|<3/4$.  Starting at such a time $N$ after all readouts have ceased,
\[
 |z_{N+r}|=|z_N|^{2^r}\longrightarrow0.
\]
The registers may have arbitrary behavior; all their values are included
in the definition of $Z_0$.
\end{proof}

For the rest of the passage analysis, assume $t_n\to+\infty$ and $z_0$ is
finite.  Every $z_n$ is then finite, because the added terms in
\eqref{eq:map} are finite.  Choose an integer $k$ with $100k>t_0$ and set
\[
 N_k=\min\{n\ge0:t_n\ge100k\},\qquad
 N_{k+1}=\min\{n\ge0:t_n\ge100(k+1)\}.
\]
These indices exist.  By~\eqref{eq:clock-increments},
$0\le t_{N_k}-100k<1$.  We call the indices
$N_k\le n<N_{k+1}$ a \emph{complete passage} and, for this fixed passage,
write
\[
 s_n=t_n-100k,\qquad
 I=\{n:N_k\le n<N_{k+1},\ \alpha(s_n)>0\}.
\]
The symbol $s_n$ is a translated clock reading, not a new time index.
Within the passage $0\le s_n<100$, and the periodic coefficients of
\eqref{eq:map} can be evaluated at $s_n$.  In particular,
$\bar\beta_\ell(t_n)=\beta_\ell(s_n)$ there.  Every later period is also a complete
passage because the lift is already assumed unbounded.

\begin{lemma}\label{lem:passage-preparation}
In every complete passage as above, all registers are reset to zero before
the first index in $I$.  The set $I$ is nonempty, has at most five elements,
and contains an index $i$ with $\alpha(s_i)=1$.  For every $i\in I$,
\begin{equation}\label{eq:preparation-radius}
 |z_i|>\tfrac12.
\end{equation}
There also exists an index $j$ in this passage with $s_j\in[76,77]$ and
$\beta_\ell(s_j)=1$ for at least one $\ell$.
\end{lemma}

\begin{proof}
Let $r$ be the first index in the passage with $s_r\ge2$.  The preceding
reading is below $2$, so $2\le s_r<3$.  At time $r$ one has $a(s_r)=0$ and
$\alpha(s_r)=0$, whence $U_\ell(s_r,z_r)=0$ and
\[
 v_{r+1,\ell}=0\qquad(1\le \ell\le q).
\]
This conclusion holds even if any earlier register was at infinity.
Until preparation starts, the inputs remain zero, and $S(0)=0$ keeps all
registers zero.

On $(8,12)$ the clock speed is exactly one.  Its readings in that interval
therefore belong to a single translate of the integers, which gives
$|I|\le5$.  The first reading at least $9$ belongs to $[9,10)$ by
\eqref{eq:clock-increments}; there $\alpha=1$.  Similarly, the first reading
at least $76$ belongs to $[76,77)$, where the plateau-covering property of
the $\beta_\ell$ supplies a coefficient equal to one.  All these crossings
occur in the passage because $s_n$ eventually reaches $100$.

Suppose now that $i\in I$ and $|z_i|\le1/2$.  As long as
$s_n\in[8,50]$ and $|z_n|\le1/2$, there is no readout, $e(z_n)=0$, and
\[
 z_{n+1}=z_n^2,\qquad s_{n+1}=s_n+b(s_n)\le50.
\]
The clock cannot decrease, and $|z_{n+1}|\le1/4\le1/2$.  These two
conditions thus persist by induction from time $i$, keeping $s_n\le50$
forever.  This contradicts the assumed complete passage.  This also
excludes the endpoint $|z_i|=1/2$ and proves~\eqref{eq:preparation-radius}.
\end{proof}

\subsection{Prediction and linearity up to a register's readout}

Fix one complete passage.  The next two lemmas identify the value of a
register at its readout.  The prediction uses the uncontrolled map $G$,
whose main coordinate need not agree with the actual main coordinate after
a readout.  Only their clock coordinates must continue to agree.

\begin{lemma}\label{lem:clock-prediction}
For a given $\ell$, at most one index $j$ in the passage satisfies
$\beta_\ell(s_j)>0$.  If such $j$ exists, then, for every $i\in I$ and every
integer $m\ge1$,
\begin{equation}\label{eq:prediction-single-hit}
 \beta_\ell\bigl(T_m(s_i,z_i)\bigr)=
 \begin{cases}
  \beta_\ell(s_j),&m=j-i,\\
  0,&m\ne j-i.
 \end{cases}
\end{equation}
Here $\beta_\ell$ is the function on $\mathbb R$ used in~\eqref{eq:input},
not its periodization.
\end{lemma}

\begin{proof}
Fix $i\in I$, and let $r_{70}$ and $r_{82}$ be the first indices after $i$
with $s_{r_{70}}\ge70$ and $s_{r_{82}}\ge82$, respectively.  Before
$r_{70}$ all readout coefficients vanish.  The actual pair $(s_n,z_n)$
therefore follows $G$ exactly from time $i$ through time $r_{70}$.
At the latter time $s_{r_{70}}\in[70,71)$.  Throughout $[70,82]$ one has
$b=1$, so both the actual and predicted clock increments are one,
independently of their respective main coordinates.  Starting with the
equal clocks at time $r_{70}$, induction gives
\begin{equation}\label{eq:clock-prediction-agreement}
 T_{n-i}(s_i,z_i)=s_n\qquad(i<n\le r_{82}),
\end{equation}
including after any main-coordinate additions.  The predicted clock is
nondecreasing at all later times.  Once it reaches $82$, it cannot return
to any support of the nonperiodic $\beta_\ell$, all of which lie in $(74,80)$.

The support of a fixed $\beta_\ell$ lies in an interval of length less than
one.  If $\beta_\ell(s_j)>0$, the next clock reading is $s_j+1$, strictly
beyond that support interval.  Monotonicity then rules out any later
positive value of this $\beta_\ell$; the same argument rules out two positive
values in either order.  Equation~\eqref{eq:clock-prediction-agreement}
identifies every possible predicted hit with an actual hit in this
passage, and later predicted hits are impossible.  This proves
\eqref{eq:prediction-single-hit}.
\end{proof}

\begin{lemma}\label{lem:register-linearity}
Let $j$ be an index in the fixed complete passage and let $\ell$ satisfy
$\beta_\ell(s_j)>0$.  Then $|v_{j,\ell}|\le1/2$, $S(v_{j,\ell})=v_{j,\ell}$, and
\begin{equation}\label{eq:register-at-readout}
 v_{j,\ell}=\frac{\beta_\ell(s_j)}{10}
 \sum_{i\in I}\alpha(s_i)\chi(z_i)
 \left(\frac{z_i}{|z_i|}\right)^{2^{j-i+1}}.
\end{equation}
This conclusion does not assume any alignment of the summands' phases.
\end{lemma}

\begin{proof}
Put $i_0=\min I$ and, for each $i\in I$, define
\[
 c_i=\frac{\beta_\ell(s_j)}{10}\alpha(s_i)\chi(z_i)
      \left(\frac{z_i}{|z_i|}\right)^{2^{j-i+1}}.
\]
The phases are defined by~\eqref{eq:preparation-radius}.
Lemma~\ref{lem:clock-prediction} and~\eqref{eq:input} give the exact input
\begin{equation}\label{eq:input-target-time}
 U_\ell(s_i,z_i)=L^{1+i-j}c_i\qquad(i\in I).
\end{equation}
The input at time $i$ first enters $v_{i+1,\ell}$; it receives $j-i-1$
subsequent factors of $L$ before $v_{j,\ell}$ is read.  Thus its coefficient
would be $L^{1+i-j}L^{j-i-1}=1$ if the register remained in the region
where $S$ is the identity.  We verify that region condition directly.

For every integer $u$ with $i_0\le u\le j$, define the candidate value
\[
 w_u=L^{u-j}\sum_{\substack{i\in I\\i<u}}c_i.
\]
The triangle inequality and $|I|\le5$ give
\begin{equation}\label{eq:register-partial-sum-bound}
 |w_u|\le\frac{|I|}{10}L^{u-j}\le\tfrac12.
\end{equation}
By Lemma~\ref{lem:passage-preparation}, $v_{i_0,\ell}=0=w_{i_0}$.
For $i_0\le u<j$ the clock lies between the preparation and readout
intervals, so $a(s_u)=L$.  The input is zero unless $u\in I$, and otherwise
it is $L^{1+u-j}c_u$ by~\eqref{eq:input-target-time}.  Consequently
\[
 w_{u+1}=Lw_u+U_\ell(s_u,z_u).
\]
If $v_{u,\ell}=w_u$, the bound~\eqref{eq:register-partial-sum-bound} implies
$S(v_{u,\ell})=w_u$, so the actual register update gives
$v_{u+1,\ell}=w_{u+1}$.  Induction proves $v_{u,\ell}=w_u$ all the way through
$u=j$.  Every preparation index precedes $j$, hence $w_j=\sum_{i\in I}c_i$,
which proves~\eqref{eq:register-at-readout} and the stated bounds.

This induction concerns only the particular register whose readout is at
$j$.  Another register can already have been read and can subsequently
leave the unit disk.  By Lemma~\ref{lem:clock-prediction}, its readout
coefficient is zero for the rest of this passage.  Registers do not enter
one another's updates, so such later nonlinear behavior has no effect on
the induction just proved.
\end{proof}

\subsection{Phase alignment and escape}

The register values have now been established without any radial or phase
assumption beyond~\eqref{eq:preparation-radius}.  Their phases can therefore
be used to prove the main-coordinate growth.

\begin{lemma}\label{lem:passage-escape}
In a complete passage with finite main coordinate, let $i_0=\min I$.
For every $i_0\le n<N_{k+1}$, $z_n\ne0$ and
\begin{equation}\label{eq:square-lower-bound}
 |z_{n+1}|\ge |z_n|^2,
 \qquad
 \frac{z_{n+1}}{|z_{n+1}|}
   =\left(\frac{z_n}{|z_n|}\right)^2.
\end{equation}
Moreover, some index $N$ with $i_0\le N<N_{k+1}$ satisfies $|z_N|>1$.
\end{lemma}

\begin{proof}
Set $\theta=z_{i_0}/|z_{i_0}|$.  Before the first readout, the main
coordinate squares exactly, so it stays nonzero and its phase at time $n$
is $\theta^{2^{n-i_0}}$.  In particular this holds at every preparation
index $i\in I$, since all preparation precedes all readouts.

Proceed chronologically through the readout indices.  Suppose the phase at
time $j$ is $\theta^{2^{j-i_0}}$ and $z_j\ne0$.  For every $i\in I$ the
phase appearing in~\eqref{eq:register-at-readout} is
\[
 \left(\frac{z_i}{|z_i|}\right)^{2^{j-i+1}}
 =\left(\theta^{2^{i-i_0}}\right)^{2^{j-i+1}}
 =\theta^{2^{j-i_0+1}},
\]
which is the phase of $z_j^2$.  Lemma~\ref{lem:register-linearity} applies
separately to each $\ell$ with $\beta_\ell(s_j)>0$.  The other channels have zero
coefficient in the main update.  Thus that update is exactly
\begin{equation}\label{eq:aligned-main-update}
 z_{j+1}=
 \left(
 |z_j|^2+
 2\sum_{\ell=1}^q\beta_\ell(s_j)^2
       \sum_{i\in I}\alpha(s_i)\chi(z_i)
 \right)\theta^{2^{j-i_0+1}}.
\end{equation}
One factor of $\beta_\ell(s_j)$ comes from the input and the other from the
readout in~\eqref{eq:map}; both are essential to this coefficient.
Every added coefficient is nonnegative, and the baseline $|z_j|^2$ is
strictly positive.  This proves nonvanishing, phase doubling, and the
square lower bound at this readout, even if all its inputs vanish.
Between readouts there is exact squaring.  Induction therefore proves
\eqref{eq:square-lower-bound} through the end of the passage, including
simultaneous readouts of several channels and successive readouts of
different channels.

If some preparation value already has $|z_i|>1$, take $N=i$.  Otherwise
$1/2<|z_i|\le1$ for every $i\in I$.  Choose $i_*$ with
$\alpha(s_{i_*})=1$, supplied by Lemma~\ref{lem:passage-preparation}.
The prescribed cutoff satisfies $\chi(z_{i_*})=1$, including when
$|z_{i_*}|=1$.  The same lemma supplies a readout $j_*$ and a channel $\ell_*$
with $s_{j_*}\in[76,77]$ and $\beta_{\ell_*}(s_{j_*})=1$.
In~\eqref{eq:aligned-main-update}, the single pair $(\ell_*,i_*)$ therefore
adds $2$ to the radius, and all other terms have the same phase.  Hence
\[
 |z_{j_*+1}|\ge |z_{j_*}|^2+2>2.
\]
Taking $N=j_*+1$ proves the claim; this index is still in the passage
because $s_{j_*+1}\le78<100$.
\end{proof}

\begin{proof}[Proof of Proposition~\ref{prop:attraction}]
If the lift is bounded, Lemma~\ref{lem:bounded-clock} proves the asserted
coordinate limits.  In a product metric the distance to $Z_0$ consequently
tends to zero, irrespective of the registers.

If $z_0=\infty$, it remains at infinity, all inputs vanish, and the clock
advances by one at every step.  A translate of the integer lattice meets
the reset interval $[2,5]$ modulo $100$.  At such a step $a=0$, so every
register becomes zero, regardless of its initial value, and stays zero.
The orbit then lies in $Z_\infty$.

It remains to consider an unbounded lift with finite $z_0$.  Choose any
complete passage and an index $N$ with $|z_N|>1$ given by
Lemma~\ref{lem:passage-escape}.  Its square lower bound continues through
the rest of that passage.  We spell out why it also holds at every later
time.  After the current readout region $(74,80)$, all readout multipliers
vanish until the next period's region $(174,180)$ in the same translated
coordinate.  In particular, the main variable squares exactly across the
period boundary and throughout the next reset and preparation intervals.
An old register, however large, cannot alter $z$ there.  It is erased by
the next reset before any new preparation input is entered.

Every later period is a complete passage by the unbounded-lift assumption.
Lemma~\ref{lem:passage-escape} therefore applies afresh from its first
preparation index through its end.  The intervening indices, before that
first preparation, have no readout and give exact squaring.  These intervals
cover every time after $N$, and hence
\begin{equation}\label{eq:escape-growth}
 |z_{N+r}|\ge |z_N|^{2^r}\qquad(r\ge0).
\end{equation}
This argument does not require $\chi$ to vanish immediately after the
radius exceeds one.  New inputs in later periods still have nonnegative
weights $\alpha\chi$, satisfy the register bound, and give the aligned
additions in~\eqref{eq:aligned-main-update}.  Thus they preserve
\eqref{eq:escape-growth} even while $\chi$ remains nonzero.

It follows that $z_n\to\infty$ in $\Sigma$.  Since $\chi$ vanishes in a
neighborhood of infinity, there is a time after which $\chi(z_n)=0$ at
every step.  All inputs $U_\ell(t_n,z_n)$ are then zero.  A later reset, which
exists because the clock lift is unbounded, makes every register zero;
the recurrence and $S(0)=0$ keep them zero forever.  The distance to
$Z_\infty$ therefore tends to zero.

We have proved the assertion in a product metric.  All compatible metrics
on the compact space $M$ are uniformly equivalent, so the same convergence
in distance holds for any such metric.  The conclusion is pointwise in
the initial state; no uniform entry time is asserted.
\end{proof}
